\begin{textsample}{POS Dim 1 – al – Score 34.00 – al/al\_em\_41.txt}  \label{ex:f1_pos_157}
COMPREENDENDO AS COMPETÊNCIAS ESPECÍFICAS DA ÁREA DE CIÊNCIAS HUMANAS E SOCIAIS APLICADAS Após a \textbf{exposição} e análise das principais contribuições das Ciências Humanas e Sociais Aplicadas na formação dos/as estudantes alagoanos e, do apontamento das principais características dos saberes escolares de cada \textbf{disciplina} formadora dessa área de conhecimento, passamos a expor e discutir as seis competências a serem desenvolvidas na prática docente.
Esse esforço justifica -se por compreendermos a importância de não deixar dúvidas quanto aos aspectos \textbf{que} envolvem as competências da área de Ciências Humanas e Sociais Aplicadas.
A ideia é auxiliar o entendimento \textbf{teórico} dos conceitos expostos em cada competência e a necessidade do desenvolvimento de práticas pedagógicas \textbf{condizentes}.
Será a partir do entendimento das competências específicas \textbf{que} o organizador curricular, a próxima parte do nosso documento, será mais bem compreendido e colocado em prática.
A partir de então seguimos à transcrição de cada \textbf{uma} das competências específicas e suas respectivas análises.
Competência 1 - “ Analisar processos políticos, econômicos, sociais, ambientais e culturais nos âmbitos local, regional, nacional e mundial em diferentes tempos, a partir da pluralidade de procedimentos epistemológicos, científicos e tecnológicos, de modo a compreender e posicionar -se criticamente em relação a eles, considerando diferentes pontos de vista e tomando decisões \textbf{baseadas} em argumentos e fontes de natureza científica. ” Essa competência trata de questões epistemológicas referentes às Ciências Humanas.
Pretende \textbf{despertar} no/a estudante o \textbf{espírito} investigativo necessário à análise social nos diferentes tempos e espaços.
A competência valoriza a investigação científica – empírica e filosófica – e estimula o desenvolvimento de análises por meio de conceitos e \textbf{fundamentos} \textbf{que} irão mobilizar e criticar ao buscar compreender e interpretar a realidade social no tempo e no espaço.
Nesse sentido, pretende -se estimular o/a estudante a reelaborar os seus saberes e conceitos prévios ao entrar em contato com as ferramentas e procedimentos \textbf{metodológicos} próprios das ciências humanas, visando o desenvolvimento de \textbf{uma} postura analítica e crítica mediante processos políticos, econômicos, sociais, ambientais e culturais. \textbf{Uma} das possibilidades de trabalho para o desenvolvimento dessa competência é o \textbf{estímulo} à investigação científica, a partir de objetos empíricos do cotidiano dos/as estudantes, mostrando \textbf{que} é possível analisar aspectos do dia a dia de diferentes formas.
Sugere -se, então, trazer os pressupostos dos/as estudantes acerca do objeto analisado para, então, elaborar junto com eles, as possíveis explicações a partir de \textbf{teorias} de análise do objeto em questão.
Pode -se pensar também no desenvolvimento de algum tipo de intervenção social relacionados ao objeto pesquisado.
Competência 2 - “ Analisar a formação de territórios e \textbf{fronteiras} em diferentes tempos e espaços, mediante a compreensão das relações de poder \textbf{que} determinam as \textbf{territorialidades} e o papel geopolítico dos Estados-nações ”.
A competência em questão reflete a capacidade do \textbf{educando} em discutir, compreender e se posicionar diante dos variados processos de formação dos territórios e das \textbf{fronteiras}, considerando os seguintes aspectos : a ) as relações de poder ; b ) o papel do Estado ; c ) as diferentes dinâmicas de produção ; d ) o tempo ; e ) o espaço ; e f ) as dimensões simbólicas.
Ao considerar as relações de poder percebe -se a existência de interesses e disputas em torno dos territórios e das definições de \textbf{fronteiras} \textbf{que} levam a sua produção, manutenção e ressignificação cultural, política, religiosa, econômica, etc.
Compreende -se os papéis políticos, econômicos e culturais do Estado e suas formas de atuação ( ou não atuação ) sobre os territórios e \textbf{fronteiras} ( de produção, manutenção, ressignificação etc. ), assim como as dinâmicas conduzidas \textbf{face} a outros territórios ( de controle, proteção, relação, interação, etc ).
Entende -se as diferentes dinâmicas ( econômicas, culturais, religiosas, política, naturais etc. ) de (re)produção das \textbf{fronteiras} e dos territórios tomando o tempo e o espaço como variáveis importantes, já \textbf{que} essas dinâmicas apresentam variações de acordo com os lugares e as épocas, sendo processos marcados por múltiplos simbolismos.
Tal competência converte -se na capacidade de pensar o conceito de poder, a formação dos Estados \textbf{modernos}, do território brasileiro e alagoano, o \textbf{que} se desdobra na \textbf{problematização} da formação das identidades, dos sentimentos de pertencimento, da alteridade, bem como dos conflitos e transformações produzidos pelas diferentes \textbf{territorialidades} \textbf{contemporâneas}, as quais se manifestam em múltiplas escalas ( local, regional, estadual, nacional e mundial ).
Competência 3 - “ Analisar e avaliar criticamente as relações de diferentes grupos, povos e sociedades com a natureza ( produção, distribuição e consumo ) e seus impactos econômicos e socioambientais, com vistas à proposição de alternativas \textbf{que} respeitem e promovam a consciência, a ética socioambiental e o consumo responsável em âmbito local, regional, nacional e global ”.
Esta competência refere -se a capacidade de analisar e avaliar criticamente as relações entre o \textbf{homem} e a natureza, levando em consideração a economia, a política e as culturas dos diferentes grupos humanos.
Entende -se o processo histórico como fundamental para a compreensão dos usos do espaço ocupado e como a sociedade e o meio ambiente estão entrelaçados.
Importa analisar \textbf{que} os impactos ambientais gerados pelo consumismo podem causar danos socioambientais e econômicos nos níveis global e local.
Esse entendimento é essencial para \textbf{que} os/as nossos/as estudantes analisem e avaliem como ocorrem os fluxos de mercadoria, bem como as ressignificações dos fatores de produção.
Da mesma forma, é necessário compreender como as populações \textbf{que} \textbf{vivem} nos países ainda em desenvolvimento sofrem danos irreparáveis causados pela exploração ambiental dos países ricos.
Competência 4 - “ Analisar as relações de produção, capital e trabalho em diferentes territórios, contextos e culturas, discutindo o papel dessas relações na construção, consolidação e transformação das sociedades ”.
Essa competência possibilita a compreensão das diferentes realidades do trabalho e seus múltiplos significados na sociedade \textbf{moderna} e \textbf{contemporânea}, do Mundo, do Brasil e, em especial, de Alagoas.
A proposta é refletir sobre a \textbf{divisão} do trabalho, da \textbf{hierarquia} e das desigualdades sociais \textbf{que} estão \textbf{imbricadas} nesse processo.
Faz -se necessário \textbf{uma} reflexão \textbf{dialética} com as dimensões históricas, geográficas, sociológicas, \textbf{antropológicas} e filosóficas, \textbf{visto} \textbf{que} a análise das relações de produção do capital e do trabalho, em diversas sociedades podem contribuir para compreender melhor as formas de exploração da mão de obra presentes na realidade e nas diferentes regiões do estado, destacando \textbf{aqui} o trabalho na cana-de-açúcar, no turismo, na pesca, no artesanato, no comércio etc.
Esta competência contribui para \textbf{que} os/as estudantes compreendam as complexas relações \textbf{que} determinam a atividade produtiva capitalista e outras formas históricas de relações de produção em diversos espaços.
Partindo da relação existente entre valor/trabalho/capital os/as estudantes se orientarão na vida prática, questionando e procurando transformar as suas realidades e da sociedade como \textbf{um} todo.
Competência 5 - “ Identificar e combater as diversas formas de injustiça, preconceito e violência, adotando princípios éticos, democráticos, inclusivos e solidários, e respeitando os Direitos Humanos ”.
As Ciências Humanas e Sociais Aplicadas, ao trabalhar com esta competência, apresentam diferentes concepções \textbf{teóricas} sobre a diversidade cultural e desigualdade social presentes nos dias de \textbf{hoje}.
É possível tratar \textbf{aqui} de temas como : diferenças e trocas culturais - problematizando o Etnocentrismo ainda tão destacado nos livros e materiais didáticos -, diferenças de gênero, \textbf{raça}, etnia, religião, relações de identidade e pertencimento ; e também de alteridade das comunidades quilombolas, indígenas e \textbf{ribeirinhas}, grupos estes \textbf{que} sempre \textbf{viveram} e ainda \textbf{vivem} à margem da sociedade alagoana.
Pode -se propor reflexões \textbf{que} busquem desnaturalizar ideias arraigadas nesses temas e \textbf{que} são reafirmadas nas experiências cotidianas de preconceitos e desigualdades sociais ; \textbf{visto} \textbf{que} estas colocam em risco os direitos humanos reafirmados pela nossa Constituição de 1988 e nas demais Leis implementadas no sistema educacional brasileiro ao longo da sua trajetória.
Neste sentido, esta competência busca fortalecer e contribuir para \textbf{que} o/a estudante desenvolva o respeito à diversidade, aos princípios éticos e aos Direitos Humanos.
Competência 6 - “ Participar do debate público de forma crítica, respeitando diferentes posições e fazendo escolhas alinhadas ao exercício da cidadania e ao seu projeto de vida, com liberdade, autonomia, consciência crítica e responsabilidade ”.
O desenvolvimento dessa competência visa o \textbf{estímulo} à participação política do/a estudante para além do exercício do voto, visando sua inserção social no debate político público.
A partir do desenvolvimento de habilidades e competências relacionadas ao aporte teórico-conceitual das ciências humanas, o/a estudante é capacitado a realizar a leitura do mundo social, político, econômico e cultural no qual está inserido.
A partir daí o/a estudante pode realizar suas escolhas e se posicionar criticamente perante sua realidade, sempre buscando sua autonomia e liberdade intelectuais, munido de \textbf{um} sentimento de tolerância e respeito às diferenças e às diversidades existentes nas sociedades.

% matched lemmas: antropológico, aqui, basear, condizente, contemporâneo, despertar, dialético, disciplina, divisão, educar, espírito, estímulo, exposição, face, fronteira, fundamento, hierarquia, hoje, homem, imbricar, metodológico, moderno, problematização, que, raça, ribeirinho, teoria, territorialidade, teórico, um, uma, ver, viver
\end{textsample}
